% !TeX program = pdflatex
% SPDX-License-Identifier: MIT
\documentclass[10pt,a4paper]{article}
\usepackage[margin=23mm]{geometry}
\usepackage{xcolor}
\usepackage{multicol}
\usepackage{devicon}
\usepackage[hidelinks]{hyperref}

\newcommand{\DeviconCatalogEntry}[1]{%
  \noindent\makebox[1.8em][l]{\devicon[artifact]{#1}}%
  {\footnotesize\ttfamily\detokenize{#1}}\par
}
\title{Devicon for \LaTeX}
\author{Johnny Myung Won Lee}
\date{Version 0.2.2 \quad $\mid$ \quad 27th November 2026}

\begin{document}
\maketitle
\noindent
Use the Devicon icon font for programming languages and development tools in
documents compiled with pdfLaTeX, XeLaTeX or LuaLaTeX. This package bundles Devicon v2.18.0 from upstream revision \href{https://github.com/devicons/devicon/commit/32fe1d3b3f36ad11f90ab94f7db59738139c2bf6}{\texttt{32fe1d3}}.

\section{Getting started}
To try the example, upload the Overleaf release ZIP through
\textbf{New Project} $\to$ \textbf{Upload Project} and compile
\texttt{example.tex} with pdfLaTeX.

\subsection{Add Devicon to an existing Overleaf project}
Extract the Overleaf ZIP or open the \texttt{devicon/} folder in the
extracted CTAN archive. Use \textbf{Upload} in Overleaf's file list and put
the following files in the project root, alongside your main \texttt{.tex} file.
\begin{itemize}
  \setlength{\itemsep}{0pt}
  \setlength{\parsep}{0pt}
  \item[$\bullet$] \texttt{devicon.sty}, \texttt{devicon-glyphs.def}, \texttt{devicon-2.18.0.ttf}
  \item[$\bullet$] \texttt{devicon-pdftex.map} and all six \verb|devicon-2-18-0-*.enc|
    and six \verb|devicon-2-18-0-*.tfm| files
\end{itemize}
Add \verb|\usepackage{devicon}| to your preamble and a \verb|\devicon{...}|
command in the document body. Click \textbf{Recompile}. For local TEXMF
installation paths, see \texttt{README.md}.

\begin{verbatim}
\documentclass{article}
\usepackage{devicon}
\begin{document}
\devicon{python-plain} Python
\end{document}
\end{verbatim}

\subsection{When will Overleaf include Devicon?}
Devicon is published in Devicon but not mounted on TeX Live distribution.
Once it is, you can load it directly once your project's TeX Live version includes it.
Standard Overleaf updates annually, and the opt-in Rolling TeX Live
experiment uses monthly builds. See the
\href{https://docs.overleaf.com/troubleshooting-and-support/tex-live}{Overleaf TeX Live documentation}.

\section{Using icons}
\verb|\devicon{python-plain}| produces \devicon{python-plain}.
Names are case-sensitive. The CSS prefix \texttt{devicon-} is optional.
Choose the font variant shown in the catalog, including
its suffix such as \texttt{-plain} or \texttt{-original}.

\medskip\noindent
Icons inherit the current text size and color. For example,
\verb|{\Large\devicon{git-plain}}| changes size and
\verb|\textcolor{red}{\devicon{git-plain}}| colors the glyph using
\texttt{xcolor}, as in \textcolor{red}{\devicon{git-plain}}.
The font provides single-color glyphs. Icons are omitted from
\texttt{hyperref} PDF bookmarks.

\clearpage
\section{Accessible icons}
Devicon glyphs use Unicode Private Use Area code points. The package supplies
PDF \texttt{ActualText} derived from the icon name so there is readable text
in their place. Use \texttt{actualtext} to supply a plain-text label suited
to the context. Use \texttt{artifact} for an icon whose meaning is already
present in the surrounding text.

\begin{verbatim}
\devicon[actualtext={GitHub}]{github-original}
\devicon[artifact]{github-original} GitHub
\end{verbatim}

When LaTeX PDF tagging is active, meaningful icons use a \texttt{Span} and
decorative icons are marked as artifacts. Otherwise, \texttt{accsupp}
provides replacement text for text extraction. Document-wide accessibility
requires a tagged document and validation of the finished PDF. See
\href{https://tagging-project.latex-project.org/documentation/usage-instructions}%
{the LaTeX accessibility guidance}.

\section{Credits and license}
The package code is released under the MIT License. The bundled Devicon font
and its source are from \href{https://github.com/devicons/devicon}{Devicon}
and retain the upstream MIT license. Product names
and logos belong to their respective owners. Follow the relevant brand policies
when using the icons. See \texttt{LICENSE} and \texttt{LICENSE-DEVICON}.

\section{Icon catalog}
The entries below are the font names accepted by \verb|\devicon|.
\begin{multicols}{2}
\setlength{\parindent}{0pt}
\setlength{\parskip}{1pt}
% Generated from Devicon v2.18.0 by scripts/release.py.
\DeviconCatalogEntry{aarch64-line}
\DeviconCatalogEntry{aarch64-line-wordmark}
\DeviconCatalogEntry{aarch64-plain}
\DeviconCatalogEntry{aarch64-plain-wordmark}
\DeviconCatalogEntry{adonisjs-original}
\DeviconCatalogEntry{adonisjs-original-wordmark}
\DeviconCatalogEntry{adonisjs-plain}
\DeviconCatalogEntry{adonisjs-plain-wordmark}
\DeviconCatalogEntry{aerospike-original}
\DeviconCatalogEntry{aerospike-original-wordmark}
\DeviconCatalogEntry{aerospike-plain}
\DeviconCatalogEntry{aerospike-plain-wordmark}
\DeviconCatalogEntry{aframe-original-wordmark}
\DeviconCatalogEntry{aframe-plain}
\DeviconCatalogEntry{aframe-plain-wordmark}
\DeviconCatalogEntry{aftereffects-plain}
\DeviconCatalogEntry{aiassistant-plain}
\DeviconCatalogEntry{aiassistant-plain-wordmark}
\DeviconCatalogEntry{akka-plain}
\DeviconCatalogEntry{akka-plain-wordmark}
\DeviconCatalogEntry{algolia-original}
\DeviconCatalogEntry{algolia-original-wordmark}
\DeviconCatalogEntry{almalinux-plain}
\DeviconCatalogEntry{almalinux-plain-wordmark}
\DeviconCatalogEntry{alpinejs-original}
\DeviconCatalogEntry{alpinejs-original-wordmark}
\DeviconCatalogEntry{alpinejs-plain}
\DeviconCatalogEntry{alpinejs-plain-wordmark}
\DeviconCatalogEntry{amazonwebservices-line}
\DeviconCatalogEntry{amazonwebservices-line-wordmark}
\DeviconCatalogEntry{amazonwebservices-plain}
\DeviconCatalogEntry{amazonwebservices-plain-wordmark}
\DeviconCatalogEntry{anaconda-original}
\DeviconCatalogEntry{anaconda-original-wordmark}
\DeviconCatalogEntry{anaconda-plain}
\DeviconCatalogEntry{anaconda-plain-wordmark}
\DeviconCatalogEntry{android-plain}
\DeviconCatalogEntry{android-plain-wordmark}
\DeviconCatalogEntry{androidstudio-plain}
\DeviconCatalogEntry{androidstudio-plain-wordmark}
\DeviconCatalogEntry{angular-plain}
\DeviconCatalogEntry{angular-plain-wordmark}
\DeviconCatalogEntry{angularjs-plain}
\DeviconCatalogEntry{angularjs-plain-wordmark}
\DeviconCatalogEntry{angularmaterial-plain}
\DeviconCatalogEntry{ansible-plain}
\DeviconCatalogEntry{ansible-plain-wordmark}
\DeviconCatalogEntry{ansys-plain}
\DeviconCatalogEntry{ansys-plain-wordmark}
\DeviconCatalogEntry{antdesign-plain}
\DeviconCatalogEntry{antdesign-plain-wordmark}
\DeviconCatalogEntry{antigravity-plain}
\DeviconCatalogEntry{apache-line}
\DeviconCatalogEntry{apache-line-wordmark}
\DeviconCatalogEntry{apache-plain}
\DeviconCatalogEntry{apache-plain-wordmark}
\DeviconCatalogEntry{apacheairflow-plain}
\DeviconCatalogEntry{apacheairflow-plain-wordmark}
\DeviconCatalogEntry{apacheignite-line}
\DeviconCatalogEntry{apacheignite-original}
\DeviconCatalogEntry{apacheignite-plain}
\DeviconCatalogEntry{apacheignite-plain-wordmark}
\DeviconCatalogEntry{apachekafka-original}
\DeviconCatalogEntry{apachekafka-original-wordmark}
\DeviconCatalogEntry{apachekafka-plain}
\DeviconCatalogEntry{apachekafka-plain-wordmark}
\DeviconCatalogEntry{apachespark-original}
\DeviconCatalogEntry{apachespark-plain}
\DeviconCatalogEntry{apachespark-plain-wordmark}
\DeviconCatalogEntry{apex-original}
\DeviconCatalogEntry{apex-plain}
\DeviconCatalogEntry{apify-plain}
\DeviconCatalogEntry{apify-plain-wordmark}
\DeviconCatalogEntry{apl-plain}
\DeviconCatalogEntry{apollographql-line}
\DeviconCatalogEntry{apollographql-line-wordmark}
\DeviconCatalogEntry{apollographql-original}
\DeviconCatalogEntry{apollographql-original-wordmark}
\DeviconCatalogEntry{apollographql-plain}
\DeviconCatalogEntry{apollographql-plain-wordmark}
\DeviconCatalogEntry{appcelerator-original}
\DeviconCatalogEntry{appcelerator-plain}
\DeviconCatalogEntry{appcelerator-plain-wordmark}
\DeviconCatalogEntry{apple-original}
\DeviconCatalogEntry{apple-plain}
\DeviconCatalogEntry{appwrite-original}
\DeviconCatalogEntry{appwrite-plain}
\DeviconCatalogEntry{appwrite-plain-wordmark}
\DeviconCatalogEntry{archlinux-plain}
\DeviconCatalogEntry{archlinux-plain-wordmark}
\DeviconCatalogEntry{arduino-plain}
\DeviconCatalogEntry{arduino-plain-wordmark}
\DeviconCatalogEntry{argocd-plain}
\DeviconCatalogEntry{argocd-plain-wordmark}
\DeviconCatalogEntry{artixlinux-plain}
\DeviconCatalogEntry{artixlinux-plain-wordmark}
\DeviconCatalogEntry{aspire-line}
\DeviconCatalogEntry{aspire-plain}
\DeviconCatalogEntry{astro-plain}
\DeviconCatalogEntry{astro-plain-wordmark}
\DeviconCatalogEntry{atom-original}
\DeviconCatalogEntry{atom-original-wordmark}
\DeviconCatalogEntry{atom-plain}
\DeviconCatalogEntry{atom-plain-wordmark}
\DeviconCatalogEntry{awk-plain-wordmark}
\DeviconCatalogEntry{axios-plain}
\DeviconCatalogEntry{axios-plain-wordmark}
\DeviconCatalogEntry{azure-plain}
\DeviconCatalogEntry{azure-plain-wordmark}
\DeviconCatalogEntry{azuredatafactory-line}
\DeviconCatalogEntry{azuredatafactory-plain}
\DeviconCatalogEntry{azuredevops-plain}
\DeviconCatalogEntry{azuresqldatabase-plain}
\DeviconCatalogEntry{babel-original}
\DeviconCatalogEntry{babel-plain}
\DeviconCatalogEntry{babylonjs-plain}
\DeviconCatalogEntry{babylonjs-plain-wordmark}
\DeviconCatalogEntry{backbonejs-plain}
\DeviconCatalogEntry{backbonejs-plain-wordmark}
\DeviconCatalogEntry{ballerina-line}
\DeviconCatalogEntry{ballerina-line-wordmark}
\DeviconCatalogEntry{ballerina-original}
\DeviconCatalogEntry{ballerina-original-wordmark}
\DeviconCatalogEntry{ballerina-plain}
\DeviconCatalogEntry{ballerina-plain-wordmark}
\DeviconCatalogEntry{bamboo-original}
\DeviconCatalogEntry{bamboo-original-wordmark}
\DeviconCatalogEntry{bamboo-plain}
\DeviconCatalogEntry{bamboo-plain-wordmark}
\DeviconCatalogEntry{bash-plain}
\DeviconCatalogEntry{bazel-plain}
\DeviconCatalogEntry{bazel-plain-wordmark}
\DeviconCatalogEntry{beats-plain}
\DeviconCatalogEntry{behance-plain}
\DeviconCatalogEntry{behance-plain-wordmark}
\DeviconCatalogEntry{bevyengine-line}
\DeviconCatalogEntry{bevyengine-line-wordmark}
\DeviconCatalogEntry{bevyengine-plain}
\DeviconCatalogEntry{bevyengine-plain-wordmark}
\DeviconCatalogEntry{biome-line}
\DeviconCatalogEntry{biome-line-wordmark}
\DeviconCatalogEntry{biome-original}
\DeviconCatalogEntry{biome-plain}
\DeviconCatalogEntry{biome-plain-wordmark}
\DeviconCatalogEntry{bitbucket-original}
\DeviconCatalogEntry{bitbucket-original-wordmark}
\DeviconCatalogEntry{bitbucket-plain}
\DeviconCatalogEntry{bitbucket-plain-wordmark}
\DeviconCatalogEntry{blazor-line}
\DeviconCatalogEntry{blazor-original}
\DeviconCatalogEntry{blazor-plain}
\DeviconCatalogEntry{blender-original}
\DeviconCatalogEntry{blender-original-wordmark}
\DeviconCatalogEntry{bootstrap-plain}
\DeviconCatalogEntry{bootstrap-plain-wordmark}
\DeviconCatalogEntry{bower-line}
\DeviconCatalogEntry{bower-line-wordmark}
\DeviconCatalogEntry{bower-plain}
\DeviconCatalogEntry{bower-plain-wordmark}
\DeviconCatalogEntry{browserstack-line}
\DeviconCatalogEntry{browserstack-line-wordmark}
\DeviconCatalogEntry{browserstack-plain}
\DeviconCatalogEntry{browserstack-plain-wordmark}
\DeviconCatalogEntry{bruno-plain}
\DeviconCatalogEntry{bruno-plain-wordmark}
\DeviconCatalogEntry{bulma-plain}
\DeviconCatalogEntry{bun-line}
\DeviconCatalogEntry{bun-plain}
\DeviconCatalogEntry{burpsuite-line}
\DeviconCatalogEntry{burpsuite-line-wordmark}
\DeviconCatalogEntry{burpsuite-plain}
\DeviconCatalogEntry{burpsuite-plain-wordmark}
\DeviconCatalogEntry{c-line}
\DeviconCatalogEntry{c-line-wordmark}
\DeviconCatalogEntry{c-original}
\DeviconCatalogEntry{c-original-wordmark}
\DeviconCatalogEntry{c-plain}
\DeviconCatalogEntry{c-plain-wordmark}
\DeviconCatalogEntry{cachyos-plain}
\DeviconCatalogEntry{cairo-plain}
\DeviconCatalogEntry{cairo-plain-wordmark}
\DeviconCatalogEntry{cake-plain}
\DeviconCatalogEntry{cakephp-plain}
\DeviconCatalogEntry{cakephp-plain-wordmark}
\DeviconCatalogEntry{canva-original}
\DeviconCatalogEntry{canva-plain}
\DeviconCatalogEntry{capacitor-plain}
\DeviconCatalogEntry{capacitor-plain-wordmark}
\DeviconCatalogEntry{carbon-original}
\DeviconCatalogEntry{carbon-plain}
\DeviconCatalogEntry{cassandra-plain}
\DeviconCatalogEntry{cassandra-plain-wordmark}
\DeviconCatalogEntry{centos-plain}
\DeviconCatalogEntry{centos-plain-wordmark}
\DeviconCatalogEntry{ceph-original}
\DeviconCatalogEntry{ceph-plain}
\DeviconCatalogEntry{ceph-plain-wordmark}
\DeviconCatalogEntry{ceylon-plain}
\DeviconCatalogEntry{ceylon-plain-wordmark}
\DeviconCatalogEntry{chakraui-plain}
\DeviconCatalogEntry{chakraui-plain-wordmark}
\DeviconCatalogEntry{chartjs-plain}
\DeviconCatalogEntry{chartjs-plain-wordmark}
\DeviconCatalogEntry{chrome-plain}
\DeviconCatalogEntry{chrome-plain-wordmark}
\DeviconCatalogEntry{circleci-plain}
\DeviconCatalogEntry{circleci-plain-wordmark}
\DeviconCatalogEntry{clarity-plain}
\DeviconCatalogEntry{clarity-plain-wordmark}
\DeviconCatalogEntry{claude-plain}
\DeviconCatalogEntry{claudecode-original}
\DeviconCatalogEntry{claudecode-plain}
\DeviconCatalogEntry{clickhouse-plain}
\DeviconCatalogEntry{clion-plain}
\DeviconCatalogEntry{clion-plain-wordmark}
\DeviconCatalogEntry{clojure-line}
\DeviconCatalogEntry{clojure-plain}
\DeviconCatalogEntry{clojurescript-plain}
\DeviconCatalogEntry{cloudflare-plain}
\DeviconCatalogEntry{cloudflare-plain-wordmark}
\DeviconCatalogEntry{cloudflareworkers-plain}
\DeviconCatalogEntry{cloudflareworkers-plain-wordmark}
\DeviconCatalogEntry{cloudrun-line}
\DeviconCatalogEntry{cloudrun-plain}
\DeviconCatalogEntry{cmake-plain}
\DeviconCatalogEntry{cmake-plain-wordmark}
\DeviconCatalogEntry{cobol-original}
\DeviconCatalogEntry{cobol-plain}
\DeviconCatalogEntry{codeac-original}
\DeviconCatalogEntry{codeberg-line}
\DeviconCatalogEntry{codeberg-line-wordmark}
\DeviconCatalogEntry{codeberg-plain}
\DeviconCatalogEntry{codeberg-plain-wordmark}
\DeviconCatalogEntry{codecov-plain}
\DeviconCatalogEntry{codeigniter-plain}
\DeviconCatalogEntry{codeigniter-plain-wordmark}
\DeviconCatalogEntry{codepen-line}
\DeviconCatalogEntry{codepen-line-wordmark}
\DeviconCatalogEntry{codepen-original}
\DeviconCatalogEntry{codepen-original-wordmark}
\DeviconCatalogEntry{codepen-plain}
\DeviconCatalogEntry{codepen-plain-wordmark}
\DeviconCatalogEntry{codex-plain}
\DeviconCatalogEntry{coffeescript-original}
\DeviconCatalogEntry{coffeescript-original-wordmark}
\DeviconCatalogEntry{coffeescript-plain}
\DeviconCatalogEntry{coffeescript-plain-wordmark}
\DeviconCatalogEntry{composer-line}
\DeviconCatalogEntry{composer-line-wordmark}
\DeviconCatalogEntry{composer-plain}
\DeviconCatalogEntry{composer-plain-wordmark}
\DeviconCatalogEntry{confluence-line}
\DeviconCatalogEntry{confluence-line-wordmark}
\DeviconCatalogEntry{confluence-plain}
\DeviconCatalogEntry{confluence-plain-wordmark}
\DeviconCatalogEntry{consul-original}
\DeviconCatalogEntry{consul-plain}
\DeviconCatalogEntry{consul-plain-wordmark}
\DeviconCatalogEntry{contao-original}
\DeviconCatalogEntry{contao-original-wordmark}
\DeviconCatalogEntry{cordova-plain}
\DeviconCatalogEntry{cordova-plain-wordmark}
\DeviconCatalogEntry{corejs-original}
\DeviconCatalogEntry{corejs-original-wordmark}
\DeviconCatalogEntry{corejs-plain}
\DeviconCatalogEntry{corejs-plain-wordmark}
\DeviconCatalogEntry{cosmosdb-plain}
\DeviconCatalogEntry{cosmosdb-plain-wordmark}
\DeviconCatalogEntry{couchbase-original}
\DeviconCatalogEntry{couchbase-plain}
\DeviconCatalogEntry{couchbase-plain-wordmark}
\DeviconCatalogEntry{couchdb-plain}
\DeviconCatalogEntry{couchdb-plain-wordmark}
\DeviconCatalogEntry{cpanel-original}
\DeviconCatalogEntry{cpanel-original-wordmark}
\DeviconCatalogEntry{cpanel-plain}
\DeviconCatalogEntry{cpanel-plain-wordmark}
\DeviconCatalogEntry{cplusplus-line}
\DeviconCatalogEntry{cplusplus-line-wordmark}
\DeviconCatalogEntry{cplusplus-plain}
\DeviconCatalogEntry{cplusplus-plain-wordmark}
\DeviconCatalogEntry{crystal-line}
\DeviconCatalogEntry{crystal-line-wordmark}
\DeviconCatalogEntry{crystal-original}
\DeviconCatalogEntry{crystal-original-wordmark}
\DeviconCatalogEntry{crystal-plain}
\DeviconCatalogEntry{crystal-plain-wordmark}
\DeviconCatalogEntry{csharp-line}
\DeviconCatalogEntry{csharp-line-wordmark}
\DeviconCatalogEntry{csharp-plain}
\DeviconCatalogEntry{csharp-plain-wordmark}
\DeviconCatalogEntry{css-plain}
\DeviconCatalogEntry{css3-plain}
\DeviconCatalogEntry{css3-plain-wordmark}
\DeviconCatalogEntry{cucumber-plain}
\DeviconCatalogEntry{cucumber-plain-wordmark}
\DeviconCatalogEntry{cursor-original-wordmark}
\DeviconCatalogEntry{cursor-plain}
\DeviconCatalogEntry{cursor-plain-wordmark}
\DeviconCatalogEntry{cypressio-line}
\DeviconCatalogEntry{cypressio-line-wordmark}
\DeviconCatalogEntry{cypressio-plain}
\DeviconCatalogEntry{cypressio-plain-wordmark}
\DeviconCatalogEntry{d3js-plain}
\DeviconCatalogEntry{dart-plain}
\DeviconCatalogEntry{dart-plain-wordmark}
\DeviconCatalogEntry{datadog-original}
\DeviconCatalogEntry{datadog-original-wordmark}
\DeviconCatalogEntry{datadog-plain}
\DeviconCatalogEntry{datadog-plain-wordmark}
\DeviconCatalogEntry{datagrip-plain}
\DeviconCatalogEntry{datagrip-plain-wordmark}
\DeviconCatalogEntry{dataspell-plain}
\DeviconCatalogEntry{dataspell-plain-wordmark}
\DeviconCatalogEntry{datatables-line}
\DeviconCatalogEntry{datatables-original}
\DeviconCatalogEntry{datatables-plain}
\DeviconCatalogEntry{dbeaver-plain}
\DeviconCatalogEntry{debian-plain}
\DeviconCatalogEntry{debian-plain-wordmark}
\DeviconCatalogEntry{delphi-plain}
\DeviconCatalogEntry{denojs-original}
\DeviconCatalogEntry{denojs-original-wordmark}
\DeviconCatalogEntry{denojs-plain}
\DeviconCatalogEntry{denojs-plain-wordmark}
\DeviconCatalogEntry{detaspace-line}
\DeviconCatalogEntry{detaspace-line-wordmark}
\DeviconCatalogEntry{devicon-line}
\DeviconCatalogEntry{devicon-line-wordmark}
\DeviconCatalogEntry{devicon-plain}
\DeviconCatalogEntry{devicon-plain-wordmark}
\DeviconCatalogEntry{digitalocean-original}
\DeviconCatalogEntry{digitalocean-original-wordmark}
\DeviconCatalogEntry{digitalocean-plain}
\DeviconCatalogEntry{digitalocean-plain-wordmark}
\DeviconCatalogEntry{discloud-original}
\DeviconCatalogEntry{discloud-plain}
\DeviconCatalogEntry{discloud-plain-wordmark}
\DeviconCatalogEntry{discordjs-plain}
\DeviconCatalogEntry{discordjs-plain-wordmark}
\DeviconCatalogEntry{django-plain}
\DeviconCatalogEntry{django-plain-wordmark}
\DeviconCatalogEntry{djangorest-line}
\DeviconCatalogEntry{djangorest-line-wordmark}
\DeviconCatalogEntry{djangorest-plain}
\DeviconCatalogEntry{djangorest-plain-wordmark}
\DeviconCatalogEntry{docker-plain}
\DeviconCatalogEntry{docker-plain-wordmark}
\DeviconCatalogEntry{doctrine-line}
\DeviconCatalogEntry{doctrine-line-wordmark}
\DeviconCatalogEntry{doctrine-plain}
\DeviconCatalogEntry{doctrine-plain-wordmark}
\DeviconCatalogEntry{dot-net-plain}
\DeviconCatalogEntry{dot-net-plain-wordmark}
\DeviconCatalogEntry{dotnetcore-plain}
\DeviconCatalogEntry{dovecot-line}
\DeviconCatalogEntry{dovecot-line-wordmark}
\DeviconCatalogEntry{dovecot-plain}
\DeviconCatalogEntry{dreamweaver-line}
\DeviconCatalogEntry{dreamweaver-plain}
\DeviconCatalogEntry{dropwizard-plain}
\DeviconCatalogEntry{drupal-plain}
\DeviconCatalogEntry{drupal-plain-wordmark}
\DeviconCatalogEntry{duckdb-plain}
\DeviconCatalogEntry{dyalog-plain}
\DeviconCatalogEntry{dynamodb-plain}
\DeviconCatalogEntry{dynatrace-line}
\DeviconCatalogEntry{dynatrace-line-wordmark}
\DeviconCatalogEntry{dynatrace-plain}
\DeviconCatalogEntry{dynatrace-plain-wordmark}
\DeviconCatalogEntry{eclipse-plain}
\DeviconCatalogEntry{eclipse-plain-wordmark}
\DeviconCatalogEntry{ecto-original}
\DeviconCatalogEntry{ecto-plain}
\DeviconCatalogEntry{ecto-plain-wordmark}
\DeviconCatalogEntry{elasticsearch-plain}
\DeviconCatalogEntry{elasticsearch-plain-wordmark}
\DeviconCatalogEntry{electron-original}
\DeviconCatalogEntry{electron-original-wordmark}
\DeviconCatalogEntry{eleventy-plain}
\DeviconCatalogEntry{elixir-plain}
\DeviconCatalogEntry{elixir-plain-wordmark}
\DeviconCatalogEntry{elm-plain}
\DeviconCatalogEntry{elm-plain-wordmark}
\DeviconCatalogEntry{emacs-original}
\DeviconCatalogEntry{emailjs-original}
\DeviconCatalogEntry{emailjs-plain}
\DeviconCatalogEntry{emailjs-plain-wordmark}
\DeviconCatalogEntry{embeddedc-plain}
\DeviconCatalogEntry{embeddedc-plain-wordmark}
\DeviconCatalogEntry{ember-original-wordmark}
\DeviconCatalogEntry{ember-plain}
\DeviconCatalogEntry{ember-plain-wordmark}
\DeviconCatalogEntry{entityframeworkcore-line}
\DeviconCatalogEntry{entityframeworkcore-plain}
\DeviconCatalogEntry{envoy-plain}
\DeviconCatalogEntry{envoy-plain-wordmark}
\DeviconCatalogEntry{erlang-plain}
\DeviconCatalogEntry{erlang-plain-wordmark}
\DeviconCatalogEntry{eslint-line}
\DeviconCatalogEntry{eslint-line-wordmark}
\DeviconCatalogEntry{eslint-plain}
\DeviconCatalogEntry{eslint-plain-wordmark}
\DeviconCatalogEntry{expo-line}
\DeviconCatalogEntry{expo-line-wordmark}
\DeviconCatalogEntry{expo-original}
\DeviconCatalogEntry{expo-original-wordmark}
\DeviconCatalogEntry{expo-plain}
\DeviconCatalogEntry{expo-plain-wordmark}
\DeviconCatalogEntry{express-original}
\DeviconCatalogEntry{express-original-wordmark}
\DeviconCatalogEntry{facebook-original}
\DeviconCatalogEntry{facebook-plain}
\DeviconCatalogEntry{fastapi-plain}
\DeviconCatalogEntry{fastapi-plain-wordmark}
\DeviconCatalogEntry{fastify-plain}
\DeviconCatalogEntry{fastify-plain-wordmark}
\DeviconCatalogEntry{faunadb-line}
\DeviconCatalogEntry{faunadb-line-wordmark}
\DeviconCatalogEntry{faunadb-original}
\DeviconCatalogEntry{faunadb-original-wordmark}
\DeviconCatalogEntry{faunadb-plain}
\DeviconCatalogEntry{faunadb-plain-wordmark}
\DeviconCatalogEntry{feathersjs-original}
\DeviconCatalogEntry{fedora-plain}
\DeviconCatalogEntry{fiber-line}
\DeviconCatalogEntry{fiber-line-wordmark}
\DeviconCatalogEntry{fiber-plain}
\DeviconCatalogEntry{fiber-plain-wordmark}
\DeviconCatalogEntry{figma-plain}
\DeviconCatalogEntry{filamentphp-original}
\DeviconCatalogEntry{filamentphp-original-wordmark}
\DeviconCatalogEntry{filamentphp-plain}
\DeviconCatalogEntry{filamentphp-plain-wordmark}
\DeviconCatalogEntry{filezilla-line}
\DeviconCatalogEntry{filezilla-line-wordmark}
\DeviconCatalogEntry{filezilla-plain}
\DeviconCatalogEntry{filezilla-plain-wordmark}
\DeviconCatalogEntry{firebase-line}
\DeviconCatalogEntry{firebase-line-wordmark}
\DeviconCatalogEntry{firebase-plain}
\DeviconCatalogEntry{firebase-plain-wordmark}
\DeviconCatalogEntry{firebird-plain}
\DeviconCatalogEntry{firefox-plain}
\DeviconCatalogEntry{firefox-plain-wordmark}
\DeviconCatalogEntry{fish-plain}
\DeviconCatalogEntry{flask-original}
\DeviconCatalogEntry{flask-original-wordmark}
\DeviconCatalogEntry{flask-plain}
\DeviconCatalogEntry{flask-plain-wordmark}
\DeviconCatalogEntry{flutter-plain}
\DeviconCatalogEntry{flyway-original}
\DeviconCatalogEntry{flyway-original-wordmark}
\DeviconCatalogEntry{flyway-plain}
\DeviconCatalogEntry{flyway-plain-wordmark}
\DeviconCatalogEntry{forgejo-line}
\DeviconCatalogEntry{forgejo-line-wordmark}
\DeviconCatalogEntry{forgejo-plain}
\DeviconCatalogEntry{forgejo-plain-wordmark}
\DeviconCatalogEntry{fortran-original}
\DeviconCatalogEntry{foundation-plain}
\DeviconCatalogEntry{foundation-plain-wordmark}
\DeviconCatalogEntry{framermotion-original}
\DeviconCatalogEntry{framermotion-original-wordmark}
\DeviconCatalogEntry{framework7-original}
\DeviconCatalogEntry{framework7-original-wordmark}
\DeviconCatalogEntry{framework7-plain}
\DeviconCatalogEntry{framework7-plain-wordmark}
\DeviconCatalogEntry{frankenphp-plain}
\DeviconCatalogEntry{frankenphp-plain-wordmark}
\DeviconCatalogEntry{fsharp-plain}
\DeviconCatalogEntry{fusion-plain}
\DeviconCatalogEntry{gardener-line}
\DeviconCatalogEntry{gardener-plain}
\DeviconCatalogEntry{gatling-line}
\DeviconCatalogEntry{gatling-line-wordmark}
\DeviconCatalogEntry{gatling-original}
\DeviconCatalogEntry{gatling-plain}
\DeviconCatalogEntry{gatling-plain-wordmark}
\DeviconCatalogEntry{gatsby-original}
\DeviconCatalogEntry{gatsby-plain}
\DeviconCatalogEntry{gatsby-plain-wordmark}
\DeviconCatalogEntry{gazebo-plain}
\DeviconCatalogEntry{gazebo-plain-wordmark}
\DeviconCatalogEntry{gcc-line}
\DeviconCatalogEntry{gcc-line-wordmark}
\DeviconCatalogEntry{gcc-plain}
\DeviconCatalogEntry{gcc-plain-wordmark}
\DeviconCatalogEntry{gentoo-line}
\DeviconCatalogEntry{gentoo-line-wordmark}
\DeviconCatalogEntry{gentoo-plain}
\DeviconCatalogEntry{gentoo-plain-wordmark}
\DeviconCatalogEntry{ghidra-line}
\DeviconCatalogEntry{ghidra-plain}
\DeviconCatalogEntry{ghost-original}
\DeviconCatalogEntry{ghost-original-wordmark}
\DeviconCatalogEntry{gimp-line}
\DeviconCatalogEntry{gimp-line-wordmark}
\DeviconCatalogEntry{gimp-plain}
\DeviconCatalogEntry{gimp-plain-wordmark}
\DeviconCatalogEntry{git-plain}
\DeviconCatalogEntry{git-plain-wordmark}
\DeviconCatalogEntry{gitbook-line}
\DeviconCatalogEntry{gitbook-line-wordmark}
\DeviconCatalogEntry{gitbook-original}
\DeviconCatalogEntry{gitbook-original-wordmark}
\DeviconCatalogEntry{gitbook-plain}
\DeviconCatalogEntry{gitbook-plain-wordmark}
\DeviconCatalogEntry{gitea-original}
\DeviconCatalogEntry{gitea-plain}
\DeviconCatalogEntry{gitea-plain-wordmark}
\DeviconCatalogEntry{github-original}
\DeviconCatalogEntry{github-original-wordmark}
\DeviconCatalogEntry{github-plain}
\DeviconCatalogEntry{githubactions-plain}
\DeviconCatalogEntry{githubactions-plain-wordmark}
\DeviconCatalogEntry{githubcodespaces-plain}
\DeviconCatalogEntry{githubcopilot-line}
\DeviconCatalogEntry{githubcopilot-line-wordmark}
\DeviconCatalogEntry{githubcopilot-plain}
\DeviconCatalogEntry{githubcopilot-plain-wordmark}
\DeviconCatalogEntry{gitkraken-original}
\DeviconCatalogEntry{gitkraken-plain}
\DeviconCatalogEntry{gitkraken-plain-wordmark}
\DeviconCatalogEntry{gitlab-plain}
\DeviconCatalogEntry{gitlab-plain-wordmark}
\DeviconCatalogEntry{gitpod-plain}
\DeviconCatalogEntry{gitpod-plain-wordmark}
\DeviconCatalogEntry{gitter-plain}
\DeviconCatalogEntry{gitter-plain-wordmark}
\DeviconCatalogEntry{gleam-plain}
\DeviconCatalogEntry{glitch-plain}
\DeviconCatalogEntry{gnuradio-original}
\DeviconCatalogEntry{gnuradio-plain}
\DeviconCatalogEntry{gnuradio-plain-wordmark}
\DeviconCatalogEntry{go-line}
\DeviconCatalogEntry{go-original-wordmark}
\DeviconCatalogEntry{go-plain}
\DeviconCatalogEntry{go-plain-wordmark}
\DeviconCatalogEntry{godot-plain}
\DeviconCatalogEntry{godot-plain-wordmark}
\DeviconCatalogEntry{goland-plain}
\DeviconCatalogEntry{goland-plain-wordmark}
\DeviconCatalogEntry{google-original}
\DeviconCatalogEntry{google-original-wordmark}
\DeviconCatalogEntry{google-plain}
\DeviconCatalogEntry{google-plain-wordmark}
\DeviconCatalogEntry{googleappscript-plain}
\DeviconCatalogEntry{googlecloud-plain}
\DeviconCatalogEntry{googlecloud-plain-wordmark}
\DeviconCatalogEntry{googlecolab-plain}
\DeviconCatalogEntry{gradle-original}
\DeviconCatalogEntry{gradle-original-wordmark}
\DeviconCatalogEntry{gradle-plain}
\DeviconCatalogEntry{gradle-plain-wordmark}
\DeviconCatalogEntry{grafana-line}
\DeviconCatalogEntry{grafana-line-wordmark}
\DeviconCatalogEntry{grafana-plain}
\DeviconCatalogEntry{grafana-plain-wordmark}
\DeviconCatalogEntry{grafanaloki-plain}
\DeviconCatalogEntry{grafanaloki-plain-wordmark}
\DeviconCatalogEntry{grails-plain}
\DeviconCatalogEntry{graphql-plain}
\DeviconCatalogEntry{graphql-plain-wordmark}
\DeviconCatalogEntry{groovy-plain}
\DeviconCatalogEntry{grpc-plain}
\DeviconCatalogEntry{grunt-line}
\DeviconCatalogEntry{grunt-line-wordmark}
\DeviconCatalogEntry{grunt-plain}
\DeviconCatalogEntry{grunt-plain-wordmark}
\DeviconCatalogEntry{gulp-plain}
\DeviconCatalogEntry{hadoop-line}
\DeviconCatalogEntry{hadoop-line-wordmark}
\DeviconCatalogEntry{hadoop-plain}
\DeviconCatalogEntry{hadoop-plain-wordmark}
\DeviconCatalogEntry{handlebars-line}
\DeviconCatalogEntry{handlebars-line-wordmark}
\DeviconCatalogEntry{handlebars-original}
\DeviconCatalogEntry{handlebars-original-wordmark}
\DeviconCatalogEntry{handlebars-plain}
\DeviconCatalogEntry{handlebars-plain-wordmark}
\DeviconCatalogEntry{haproxy-plain}
\DeviconCatalogEntry{haproxy-plain-wordmark}
\DeviconCatalogEntry{harbor-line}
\DeviconCatalogEntry{harbor-line-wordmark}
\DeviconCatalogEntry{harbor-plain}
\DeviconCatalogEntry{harbor-plain-wordmark}
\DeviconCatalogEntry{hardhat-plain}
\DeviconCatalogEntry{hardhat-plain-wordmark}
\DeviconCatalogEntry{harvester-original}
\DeviconCatalogEntry{harvester-plain}
\DeviconCatalogEntry{harvester-plain-wordmark}
\DeviconCatalogEntry{haskell-plain}
\DeviconCatalogEntry{haskell-plain-wordmark}
\DeviconCatalogEntry{haxe-plain}
\DeviconCatalogEntry{helix-line}
\DeviconCatalogEntry{helix-plain}
\DeviconCatalogEntry{helm-line}
\DeviconCatalogEntry{helm-original}
\DeviconCatalogEntry{helm-plain}
\DeviconCatalogEntry{heroku-line}
\DeviconCatalogEntry{heroku-line-wordmark}
\DeviconCatalogEntry{heroku-original}
\DeviconCatalogEntry{heroku-original-wordmark}
\DeviconCatalogEntry{heroku-plain}
\DeviconCatalogEntry{heroku-plain-wordmark}
\DeviconCatalogEntry{hibernate-plain}
\DeviconCatalogEntry{hibernate-plain-wordmark}
\DeviconCatalogEntry{homebrew-line}
\DeviconCatalogEntry{homebrew-line-wordmark}
\DeviconCatalogEntry{homebrew-plain}
\DeviconCatalogEntry{homebrew-plain-wordmark}
\DeviconCatalogEntry{hoppscotch-plain}
\DeviconCatalogEntry{html5-plain}
\DeviconCatalogEntry{html5-plain-wordmark}
\DeviconCatalogEntry{htmx-line}
\DeviconCatalogEntry{htmx-line-wordmark}
\DeviconCatalogEntry{htmx-plain}
\DeviconCatalogEntry{htmx-plain-wordmark}
\DeviconCatalogEntry{huggingface-line}
\DeviconCatalogEntry{huggingface-line-wordmark}
\DeviconCatalogEntry{huggingface-plain}
\DeviconCatalogEntry{huggingface-plain-wordmark}
\DeviconCatalogEntry{hugo-plain}
\DeviconCatalogEntry{hugo-plain-wordmark}
\DeviconCatalogEntry{hyperv-original-wordmark}
\DeviconCatalogEntry{hyperv-plain}
\DeviconCatalogEntry{hyperv-plain-wordmark}
\DeviconCatalogEntry{hyprland-plain}
\DeviconCatalogEntry{hyprland-plain-wordmark}
\DeviconCatalogEntry{icinga-original}
\DeviconCatalogEntry{icinga-original-wordmark}
\DeviconCatalogEntry{idris-original}
\DeviconCatalogEntry{idris-plain}
\DeviconCatalogEntry{ie10-original}
\DeviconCatalogEntry{ie10-plain}
\DeviconCatalogEntry{ifttt-original}
\DeviconCatalogEntry{ifttt-plain}
\DeviconCatalogEntry{illustrator-line}
\DeviconCatalogEntry{illustrator-plain}
\DeviconCatalogEntry{inertiajs-plain}
\DeviconCatalogEntry{inertiajs-plain-wordmark}
\DeviconCatalogEntry{influxdb-original}
\DeviconCatalogEntry{influxdb-original-wordmark}
\DeviconCatalogEntry{influxdb-plain}
\DeviconCatalogEntry{influxdb-plain-wordmark}
\DeviconCatalogEntry{inkscape-plain}
\DeviconCatalogEntry{inkscape-plain-wordmark}
\DeviconCatalogEntry{insomnia-plain}
\DeviconCatalogEntry{insomnia-plain-wordmark}
\DeviconCatalogEntry{intellij-plain}
\DeviconCatalogEntry{intellij-plain-wordmark}
\DeviconCatalogEntry{ionic-original}
\DeviconCatalogEntry{ionic-original-wordmark}
\DeviconCatalogEntry{istio-original}
\DeviconCatalogEntry{istio-original-wordmark}
\DeviconCatalogEntry{jaegertracing-plain}
\DeviconCatalogEntry{jaegertracing-plain-wordmark}
\DeviconCatalogEntry{jamstack-original}
\DeviconCatalogEntry{jamstack-plain}
\DeviconCatalogEntry{jamstack-plain-wordmark}
\DeviconCatalogEntry{jasmine-original}
\DeviconCatalogEntry{jasmine-original-wordmark}
\DeviconCatalogEntry{jasmine-plain}
\DeviconCatalogEntry{jasmine-plain-wordmark}
\DeviconCatalogEntry{java-plain}
\DeviconCatalogEntry{java-plain-wordmark}
\DeviconCatalogEntry{javascript-plain}
\DeviconCatalogEntry{jeet-plain}
\DeviconCatalogEntry{jeet-plain-wordmark}
\DeviconCatalogEntry{jekyll-plain}
\DeviconCatalogEntry{jekyll-plain-wordmark}
\DeviconCatalogEntry{jenkins-line}
\DeviconCatalogEntry{jenkins-plain}
\DeviconCatalogEntry{jest-plain}
\DeviconCatalogEntry{jetbrains-plain}
\DeviconCatalogEntry{jetpackcompose-line}
\DeviconCatalogEntry{jetpackcompose-line-wordmark}
\DeviconCatalogEntry{jetpackcompose-plain}
\DeviconCatalogEntry{jetpackcompose-plain-wordmark}
\DeviconCatalogEntry{jhipster-plain}
\DeviconCatalogEntry{jhipster-plain-wordmark}
\DeviconCatalogEntry{jira-plain}
\DeviconCatalogEntry{jira-plain-wordmark}
\DeviconCatalogEntry{jiraalign-plain}
\DeviconCatalogEntry{jiraalign-plain-wordmark}
\DeviconCatalogEntry{jquery-plain}
\DeviconCatalogEntry{jquery-plain-wordmark}
\DeviconCatalogEntry{json-plain}
\DeviconCatalogEntry{jule-original}
\DeviconCatalogEntry{jule-original-wordmark}
\DeviconCatalogEntry{jule-plain}
\DeviconCatalogEntry{jule-plain-wordmark}
\DeviconCatalogEntry{julia-plain}
\DeviconCatalogEntry{julia-plain-wordmark}
\DeviconCatalogEntry{junie-original}
\DeviconCatalogEntry{junie-plain}
\DeviconCatalogEntry{junie-plain-wordmark}
\DeviconCatalogEntry{junit-line}
\DeviconCatalogEntry{junit-line-wordmark}
\DeviconCatalogEntry{junit-plain}
\DeviconCatalogEntry{junit-plain-wordmark}
\DeviconCatalogEntry{jupyter-plain}
\DeviconCatalogEntry{jupyter-plain-wordmark}
\DeviconCatalogEntry{jwt-plain}
\DeviconCatalogEntry{jwt-plain-wordmark}
\DeviconCatalogEntry{k3os-line}
\DeviconCatalogEntry{k3os-line-wordmark}
\DeviconCatalogEntry{k3os-original}
\DeviconCatalogEntry{k3os-plain}
\DeviconCatalogEntry{k3os-plain-wordmark}
\DeviconCatalogEntry{k3s-original}
\DeviconCatalogEntry{k3s-plain}
\DeviconCatalogEntry{k3s-plain-wordmark}
\DeviconCatalogEntry{k6-original}
\DeviconCatalogEntry{k6-plain}
\DeviconCatalogEntry{kaggle-original}
\DeviconCatalogEntry{kaggle-original-wordmark}
\DeviconCatalogEntry{kaggle-plain}
\DeviconCatalogEntry{kaggle-plain-wordmark}
\DeviconCatalogEntry{kaldi-line}
\DeviconCatalogEntry{kaldi-line-wordmark}
\DeviconCatalogEntry{kaldi-plain}
\DeviconCatalogEntry{kaldi-plain-wordmark}
\DeviconCatalogEntry{kalilinux-line}
\DeviconCatalogEntry{kalilinux-line-wordmark}
\DeviconCatalogEntry{kalilinux-original}
\DeviconCatalogEntry{kalilinux-plain}
\DeviconCatalogEntry{kalilinux-plain-wordmark}
\DeviconCatalogEntry{karatelabs-plain}
\DeviconCatalogEntry{karatelabs-plain-wordmark}
\DeviconCatalogEntry{karma-plain}
\DeviconCatalogEntry{kdeneon-plain}
\DeviconCatalogEntry{keras-line}
\DeviconCatalogEntry{keras-line-wordmark}
\DeviconCatalogEntry{keras-plain}
\DeviconCatalogEntry{keras-plain-wordmark}
\DeviconCatalogEntry{kibana-plain}
\DeviconCatalogEntry{kibana-plain-wordmark}
\DeviconCatalogEntry{knexjs-original}
\DeviconCatalogEntry{knexjs-plain}
\DeviconCatalogEntry{knexjs-plain-wordmark}
\DeviconCatalogEntry{knockout-plain-wordmark}
\DeviconCatalogEntry{kotlin-plain}
\DeviconCatalogEntry{kotlin-plain-wordmark}
\DeviconCatalogEntry{krakenjs-plain}
\DeviconCatalogEntry{krakenjs-plain-wordmark}
\DeviconCatalogEntry{ktor-plain}
\DeviconCatalogEntry{ktor-plain-wordmark}
\DeviconCatalogEntry{kubeflow-line}
\DeviconCatalogEntry{kubeflow-line-wordmark}
\DeviconCatalogEntry{kubeflow-plain}
\DeviconCatalogEntry{kubeflow-plain-wordmark}
\DeviconCatalogEntry{kubernetes-line}
\DeviconCatalogEntry{kubernetes-line-wordmark}
\DeviconCatalogEntry{kubernetes-plain}
\DeviconCatalogEntry{kubernetes-plain-wordmark}
\DeviconCatalogEntry{labview-plain}
\DeviconCatalogEntry{labview-plain-wordmark}
\DeviconCatalogEntry{laminas-line}
\DeviconCatalogEntry{laminas-line-wordmark}
\DeviconCatalogEntry{laminas-original}
\DeviconCatalogEntry{laminas-original-wordmark}
\DeviconCatalogEntry{laminas-plain}
\DeviconCatalogEntry{laminas-plain-wordmark}
\DeviconCatalogEntry{laravel-line}
\DeviconCatalogEntry{laravel-line-wordmark}
\DeviconCatalogEntry{laravel-original}
\DeviconCatalogEntry{laravel-original-wordmark}
\DeviconCatalogEntry{laravel-plain}
\DeviconCatalogEntry{laravel-plain-wordmark}
\DeviconCatalogEntry{laraveljetstream-original}
\DeviconCatalogEntry{laraveljetstream-plain}
\DeviconCatalogEntry{laraveljetstream-plain-wordmark}
\DeviconCatalogEntry{latex-original}
\DeviconCatalogEntry{latex-plain}
\DeviconCatalogEntry{leetcode-line}
\DeviconCatalogEntry{leetcode-line-wordmark}
\DeviconCatalogEntry{leetcode-plain}
\DeviconCatalogEntry{leetcode-plain-wordmark}
\DeviconCatalogEntry{less-plain-wordmark}
\DeviconCatalogEntry{libgdx-line}
\DeviconCatalogEntry{libgdx-line-wordmark}
\DeviconCatalogEntry{libgdx-plain}
\DeviconCatalogEntry{libgdx-plain-wordmark}
\DeviconCatalogEntry{linear-plain}
\DeviconCatalogEntry{linear-plain-wordmark}
\DeviconCatalogEntry{linkedin-plain}
\DeviconCatalogEntry{linkedin-plain-wordmark}
\DeviconCatalogEntry{linux-plain}
\DeviconCatalogEntry{linuxmint-plain}
\DeviconCatalogEntry{linuxmint-plain-wordmark}
\DeviconCatalogEntry{liquibase-original}
\DeviconCatalogEntry{liquibase-original-wordmark}
\DeviconCatalogEntry{liquibase-plain}
\DeviconCatalogEntry{liquibase-plain-wordmark}
\DeviconCatalogEntry{livewire-plain}
\DeviconCatalogEntry{livewire-plain-wordmark}
\DeviconCatalogEntry{llvm-line}
\DeviconCatalogEntry{llvm-plain}
\DeviconCatalogEntry{lodash-plain}
\DeviconCatalogEntry{logstash-plain}
\DeviconCatalogEntry{logstash-plain-wordmark}
\DeviconCatalogEntry{lovable-plain}
\DeviconCatalogEntry{lovable-plain-wordmark}
\DeviconCatalogEntry{love2d-line}
\DeviconCatalogEntry{love2d-plain}
\DeviconCatalogEntry{lua-line}
\DeviconCatalogEntry{lua-line-wordmark}
\DeviconCatalogEntry{lua-plain}
\DeviconCatalogEntry{lua-plain-wordmark}
\DeviconCatalogEntry{luau-line}
\DeviconCatalogEntry{luau-line-wordmark}
\DeviconCatalogEntry{luau-plain}
\DeviconCatalogEntry{luau-plain-wordmark}
\DeviconCatalogEntry{lumen-original}
\DeviconCatalogEntry{lumen-plain}
\DeviconCatalogEntry{magento-line}
\DeviconCatalogEntry{magento-line-wordmark}
\DeviconCatalogEntry{magento-original}
\DeviconCatalogEntry{magento-plain}
\DeviconCatalogEntry{magento-plain-wordmark}
\DeviconCatalogEntry{mapbox-original}
\DeviconCatalogEntry{mapbox-original-wordmark}
\DeviconCatalogEntry{mapbox-plain}
\DeviconCatalogEntry{mapbox-plain-wordmark}
\DeviconCatalogEntry{mariadb-original}
\DeviconCatalogEntry{mariadb-original-wordmark}
\DeviconCatalogEntry{mariadb-plain}
\DeviconCatalogEntry{mariadb-plain-wordmark}
\DeviconCatalogEntry{markdown-original}
\DeviconCatalogEntry{markdown-plain}
\DeviconCatalogEntry{materializecss-plain}
\DeviconCatalogEntry{materialui-plain}
\DeviconCatalogEntry{matlab-line}
\DeviconCatalogEntry{matlab-plain}
\DeviconCatalogEntry{matplotlib-plain}
\DeviconCatalogEntry{matplotlib-plain-wordmark}
\DeviconCatalogEntry{mattermost-original}
\DeviconCatalogEntry{mattermost-original-wordmark}
\DeviconCatalogEntry{mattermost-plain}
\DeviconCatalogEntry{mattermost-plain-wordmark}
\DeviconCatalogEntry{maven-plain}
\DeviconCatalogEntry{maven-plain-wordmark}
\DeviconCatalogEntry{maya-plain}
\DeviconCatalogEntry{maya-plain-wordmark}
\DeviconCatalogEntry{memcached-line}
\DeviconCatalogEntry{memcached-line-wordmark}
\DeviconCatalogEntry{memcached-plain}
\DeviconCatalogEntry{memcached-plain-wordmark}
\DeviconCatalogEntry{mercurial-plain}
\DeviconCatalogEntry{mercurial-plain-wordmark}
\DeviconCatalogEntry{metabase-plain}
\DeviconCatalogEntry{meteor-plain}
\DeviconCatalogEntry{meteor-plain-wordmark}
\DeviconCatalogEntry{microsoftsqlserver-line}
\DeviconCatalogEntry{microsoftsqlserver-line-wordmark}
\DeviconCatalogEntry{microsoftsqlserver-plain}
\DeviconCatalogEntry{microsoftsqlserver-plain-wordmark}
\DeviconCatalogEntry{minitab-plain}
\DeviconCatalogEntry{mithril-line}
\DeviconCatalogEntry{mithril-original}
\DeviconCatalogEntry{mithril-plain}
\DeviconCatalogEntry{mobx-plain}
\DeviconCatalogEntry{mocha-plain}
\DeviconCatalogEntry{modx-plain}
\DeviconCatalogEntry{modx-plain-wordmark}
\DeviconCatalogEntry{moleculer-original}
\DeviconCatalogEntry{moleculer-original-wordmark}
\DeviconCatalogEntry{moleculer-plain}
\DeviconCatalogEntry{moleculer-plain-wordmark}
\DeviconCatalogEntry{mongock-plain}
\DeviconCatalogEntry{mongodb-plain}
\DeviconCatalogEntry{mongodb-plain-wordmark}
\DeviconCatalogEntry{mongoose-original}
\DeviconCatalogEntry{mongoose-original-wordmark}
\DeviconCatalogEntry{monogame-line}
\DeviconCatalogEntry{monogame-line-wordmark}
\DeviconCatalogEntry{monogame-original}
\DeviconCatalogEntry{monogame-plain}
\DeviconCatalogEntry{moodle-plain}
\DeviconCatalogEntry{moodle-plain-wordmark}
\DeviconCatalogEntry{motion-original}
\DeviconCatalogEntry{motion-plain}
\DeviconCatalogEntry{msbuild-line}
\DeviconCatalogEntry{msbuild-plain}
\DeviconCatalogEntry{msdos-line}
\DeviconCatalogEntry{msdos-plain}
\DeviconCatalogEntry{mysql-original}
\DeviconCatalogEntry{mysql-plain}
\DeviconCatalogEntry{mysql-plain-wordmark}
\DeviconCatalogEntry{n8n-plain-wordmark}
\DeviconCatalogEntry{nano-plain}
\DeviconCatalogEntry{nano-plain-wordmark}
\DeviconCatalogEntry{nasm-plain}
\DeviconCatalogEntry{nasm-plain-wordmark}
\DeviconCatalogEntry{nats-plain}
\DeviconCatalogEntry{neo4j-plain}
\DeviconCatalogEntry{neo4j-plain-wordmark}
\DeviconCatalogEntry{neovim-line}
\DeviconCatalogEntry{neovim-line-wordmark}
\DeviconCatalogEntry{neovim-plain}
\DeviconCatalogEntry{neovim-plain-wordmark}
\DeviconCatalogEntry{nestjs-line}
\DeviconCatalogEntry{nestjs-line-wordmark}
\DeviconCatalogEntry{nestjs-original}
\DeviconCatalogEntry{nestjs-original-wordmark}
\DeviconCatalogEntry{nestjs-plain}
\DeviconCatalogEntry{nestjs-plain-wordmark}
\DeviconCatalogEntry{netbeans-plain}
\DeviconCatalogEntry{netbeans-plain-wordmark}
\DeviconCatalogEntry{netbox-line}
\DeviconCatalogEntry{netbox-line-wordmark}
\DeviconCatalogEntry{netbox-plain}
\DeviconCatalogEntry{netbox-plain-wordmark}
\DeviconCatalogEntry{netlify-plain}
\DeviconCatalogEntry{netlify-plain-wordmark}
\DeviconCatalogEntry{networkx-line}
\DeviconCatalogEntry{networkx-line-wordmark}
\DeviconCatalogEntry{networkx-plain}
\DeviconCatalogEntry{networkx-plain-wordmark}
\DeviconCatalogEntry{newrelic-line}
\DeviconCatalogEntry{newrelic-line-wordmark}
\DeviconCatalogEntry{newrelic-plain}
\DeviconCatalogEntry{newrelic-plain-wordmark}
\DeviconCatalogEntry{nextjs-line}
\DeviconCatalogEntry{nextjs-line-wordmark}
\DeviconCatalogEntry{nextjs-original-wordmark}
\DeviconCatalogEntry{nextjs-plain}
\DeviconCatalogEntry{nextjs-plain-wordmark}
\DeviconCatalogEntry{nginx-original}
\DeviconCatalogEntry{nginx-original-wordmark}
\DeviconCatalogEntry{nginx-plain}
\DeviconCatalogEntry{nginx-plain-wordmark}
\DeviconCatalogEntry{ngrok-line}
\DeviconCatalogEntry{ngrok-line-wordmark}
\DeviconCatalogEntry{ngrok-original}
\DeviconCatalogEntry{ngrok-original-wordmark}
\DeviconCatalogEntry{ngrok-plain}
\DeviconCatalogEntry{ngrok-plain-wordmark}
\DeviconCatalogEntry{ngrx-plain}
\DeviconCatalogEntry{nhibernate-line}
\DeviconCatalogEntry{nhibernate-line-wordmark}
\DeviconCatalogEntry{nhibernate-plain}
\DeviconCatalogEntry{nhibernate-plain-wordmark}
\DeviconCatalogEntry{nim-line}
\DeviconCatalogEntry{nim-line-wordmark}
\DeviconCatalogEntry{nim-plain}
\DeviconCatalogEntry{nim-plain-wordmark}
\DeviconCatalogEntry{nimble-plain}
\DeviconCatalogEntry{nixos-plain}
\DeviconCatalogEntry{nixos-plain-wordmark}
\DeviconCatalogEntry{nodejs-line}
\DeviconCatalogEntry{nodejs-line-wordmark}
\DeviconCatalogEntry{nodejs-plain}
\DeviconCatalogEntry{nodejs-plain-wordmark}
\DeviconCatalogEntry{nodemon-line}
\DeviconCatalogEntry{nodemon-plain}
\DeviconCatalogEntry{nodered-line}
\DeviconCatalogEntry{nodered-plain}
\DeviconCatalogEntry{nodewebkit-line}
\DeviconCatalogEntry{nodewebkit-line-wordmark}
\DeviconCatalogEntry{nodewebkit-plain}
\DeviconCatalogEntry{nodewebkit-plain-wordmark}
\DeviconCatalogEntry{nomad-original}
\DeviconCatalogEntry{nomad-plain}
\DeviconCatalogEntry{nomad-plain-wordmark}
\DeviconCatalogEntry{norg-plain}
\DeviconCatalogEntry{notion-line}
\DeviconCatalogEntry{notion-plain}
\DeviconCatalogEntry{novelwriter-plain}
\DeviconCatalogEntry{npm-original-wordmark}
\DeviconCatalogEntry{npm-plain}
\DeviconCatalogEntry{npss-plain}
\DeviconCatalogEntry{nuget-original}
\DeviconCatalogEntry{nuget-original-wordmark}
\DeviconCatalogEntry{numpy-line}
\DeviconCatalogEntry{numpy-line-wordmark}
\DeviconCatalogEntry{numpy-plain}
\DeviconCatalogEntry{numpy-plain-wordmark}
\DeviconCatalogEntry{nuxt-original}
\DeviconCatalogEntry{nuxt-plain}
\DeviconCatalogEntry{nuxt-plain-wordmark}
\DeviconCatalogEntry{nuxtjs-plain}
\DeviconCatalogEntry{nuxtjs-plain-wordmark}
\DeviconCatalogEntry{oauth-plain}
\DeviconCatalogEntry{oauth-plain-wordmark}
\DeviconCatalogEntry{objectivec-plain}
\DeviconCatalogEntry{ocaml-plain}
\DeviconCatalogEntry{ocaml-plain-wordmark}
\DeviconCatalogEntry{odin-plain}
\DeviconCatalogEntry{odin-plain-wordmark}
\DeviconCatalogEntry{ohmyzsh-plain}
\DeviconCatalogEntry{okta-plain}
\DeviconCatalogEntry{okta-plain-wordmark}
\DeviconCatalogEntry{ollama-original}
\DeviconCatalogEntry{ollama-plain}
\DeviconCatalogEntry{openal-plain}
\DeviconCatalogEntry{openapi-line}
\DeviconCatalogEntry{openapi-line-wordmark}
\DeviconCatalogEntry{openapi-plain}
\DeviconCatalogEntry{openapi-plain-wordmark}
\DeviconCatalogEntry{opencl-line}
\DeviconCatalogEntry{opencl-line-wordmark}
\DeviconCatalogEntry{opencl-plain}
\DeviconCatalogEntry{opencl-plain-wordmark}
\DeviconCatalogEntry{opencv-plain}
\DeviconCatalogEntry{opencv-plain-wordmark}
\DeviconCatalogEntry{opengl-plain}
\DeviconCatalogEntry{openstack-original}
\DeviconCatalogEntry{openstack-plain}
\DeviconCatalogEntry{openstack-plain-wordmark}
\DeviconCatalogEntry{opensuse-original}
\DeviconCatalogEntry{opensuse-original-wordmark}
\DeviconCatalogEntry{opensuse-plain}
\DeviconCatalogEntry{opensuse-plain-wordmark}
\DeviconCatalogEntry{opentelemetry-plain}
\DeviconCatalogEntry{opentelemetry-plain-wordmark}
\DeviconCatalogEntry{opera-plain}
\DeviconCatalogEntry{opera-plain-wordmark}
\DeviconCatalogEntry{opnsense-line}
\DeviconCatalogEntry{opnsense-original}
\DeviconCatalogEntry{opnsense-plain}
\DeviconCatalogEntry{opnsense-plain-wordmark}
\DeviconCatalogEntry{oracle-original}
\DeviconCatalogEntry{oracle-plain}
\DeviconCatalogEntry{ory-original}
\DeviconCatalogEntry{ory-original-wordmark}
\DeviconCatalogEntry{ory-plain}
\DeviconCatalogEntry{ory-plain-wordmark}
\DeviconCatalogEntry{p5js-original}
\DeviconCatalogEntry{p5js-original-wordmark}
\DeviconCatalogEntry{p5js-plain}
\DeviconCatalogEntry{p5js-plain-wordmark}
\DeviconCatalogEntry{packer-line}
\DeviconCatalogEntry{packer-line-wordmark}
\DeviconCatalogEntry{packer-plain}
\DeviconCatalogEntry{packer-plain-wordmark}
\DeviconCatalogEntry{pandas-line}
\DeviconCatalogEntry{pandas-line-wordmark}
\DeviconCatalogEntry{pandas-plain}
\DeviconCatalogEntry{pandas-plain-wordmark}
\DeviconCatalogEntry{passport-original-wordmark}
\DeviconCatalogEntry{passport-plain}
\DeviconCatalogEntry{passport-plain-wordmark}
\DeviconCatalogEntry{payloadcms-original}
\DeviconCatalogEntry{payloadcms-original-wordmark}
\DeviconCatalogEntry{payloadcms-plain}
\DeviconCatalogEntry{payloadcms-plain-wordmark}
\DeviconCatalogEntry{perl-plain}
\DeviconCatalogEntry{pfsense-original}
\DeviconCatalogEntry{pfsense-original-wordmark}
\DeviconCatalogEntry{pfsense-plain}
\DeviconCatalogEntry{pfsense-plain-wordmark}
\DeviconCatalogEntry{phalcon-plain}
\DeviconCatalogEntry{phoenix-original}
\DeviconCatalogEntry{phoenix-plain}
\DeviconCatalogEntry{phoenix-plain-wordmark}
\DeviconCatalogEntry{photonengine-plain}
\DeviconCatalogEntry{photonengine-plain-wordmark}
\DeviconCatalogEntry{photoshop-line}
\DeviconCatalogEntry{photoshop-plain}
\DeviconCatalogEntry{php-plain}
\DeviconCatalogEntry{phpstorm-plain}
\DeviconCatalogEntry{phpstorm-plain-wordmark}
\DeviconCatalogEntry{pixijs-plain}
\DeviconCatalogEntry{pixijs-plain-wordmark}
\DeviconCatalogEntry{playwright-plain}
\DeviconCatalogEntry{plotly-plain}
\DeviconCatalogEntry{plotly-plain-wordmark}
\DeviconCatalogEntry{pm2-line}
\DeviconCatalogEntry{pm2-line-wordmark}
\DeviconCatalogEntry{pm2-plain}
\DeviconCatalogEntry{pm2-plain-wordmark}
\DeviconCatalogEntry{pnpm-plain}
\DeviconCatalogEntry{pnpm-plain-wordmark}
\DeviconCatalogEntry{podman-plain}
\DeviconCatalogEntry{podman-plain-wordmark}
\DeviconCatalogEntry{poetry-plain}
\DeviconCatalogEntry{polygon-plain}
\DeviconCatalogEntry{polygon-plain-wordmark}
\DeviconCatalogEntry{portainer-original}
\DeviconCatalogEntry{portainer-original-wordmark}
\DeviconCatalogEntry{portainer-plain}
\DeviconCatalogEntry{portainer-plain-wordmark}
\DeviconCatalogEntry{positron-line}
\DeviconCatalogEntry{positron-plain}
\DeviconCatalogEntry{postcss-original}
\DeviconCatalogEntry{postcss-plain}
\DeviconCatalogEntry{postcss-plain-wordmark}
\DeviconCatalogEntry{postgresql-plain}
\DeviconCatalogEntry{postgresql-plain-wordmark}
\DeviconCatalogEntry{postman-plain}
\DeviconCatalogEntry{postman-plain-wordmark}
\DeviconCatalogEntry{powershell-plain}
\DeviconCatalogEntry{premierepro-plain}
\DeviconCatalogEntry{prestashop-plain}
\DeviconCatalogEntry{prestashop-plain-wordmark}
\DeviconCatalogEntry{primeng-plain}
\DeviconCatalogEntry{prisma-original}
\DeviconCatalogEntry{prisma-original-wordmark}
\DeviconCatalogEntry{prisma-plain}
\DeviconCatalogEntry{prisma-plain-wordmark}
\DeviconCatalogEntry{processing-line}
\DeviconCatalogEntry{processing-plain}
\DeviconCatalogEntry{processing-plain-wordmark}
\DeviconCatalogEntry{processwire-original}
\DeviconCatalogEntry{processwire-plain}
\DeviconCatalogEntry{processwire-plain-wordmark}
\DeviconCatalogEntry{prolog-plain}
\DeviconCatalogEntry{prolog-plain-wordmark}
\DeviconCatalogEntry{prometheus-line}
\DeviconCatalogEntry{prometheus-line-wordmark}
\DeviconCatalogEntry{prometheus-original}
\DeviconCatalogEntry{prometheus-plain}
\DeviconCatalogEntry{prometheus-plain-wordmark}
\DeviconCatalogEntry{protractor-line}
\DeviconCatalogEntry{protractor-line-wordmark}
\DeviconCatalogEntry{protractor-plain}
\DeviconCatalogEntry{protractor-plain-wordmark}
\DeviconCatalogEntry{proxmox-plain}
\DeviconCatalogEntry{proxmox-plain-wordmark}
\DeviconCatalogEntry{pug-line}
\DeviconCatalogEntry{pug-plain}
\DeviconCatalogEntry{pulsar-original}
\DeviconCatalogEntry{pulsar-original-wordmark}
\DeviconCatalogEntry{pulsar-plain}
\DeviconCatalogEntry{pulsar-plain-wordmark}
\DeviconCatalogEntry{pulumi-plain}
\DeviconCatalogEntry{pulumi-plain-wordmark}
\DeviconCatalogEntry{puppeteer-plain}
\DeviconCatalogEntry{puppeteer-plain-wordmark}
\DeviconCatalogEntry{purescript-original}
\DeviconCatalogEntry{purescript-original-wordmark}
\DeviconCatalogEntry{putty-plain}
\DeviconCatalogEntry{pycharm-plain}
\DeviconCatalogEntry{pycharm-plain-wordmark}
\DeviconCatalogEntry{pypi-plain}
\DeviconCatalogEntry{pypi-plain-wordmark}
\DeviconCatalogEntry{pyscript-plain-wordmark}
\DeviconCatalogEntry{pytest-plain}
\DeviconCatalogEntry{pytest-plain-wordmark}
\DeviconCatalogEntry{python-plain}
\DeviconCatalogEntry{python-plain-wordmark}
\DeviconCatalogEntry{pytorch-original}
\DeviconCatalogEntry{pytorch-plain}
\DeviconCatalogEntry{pytorch-plain-wordmark}
\DeviconCatalogEntry{qiskit-original}
\DeviconCatalogEntry{qiskit-plain}
\DeviconCatalogEntry{qodana-plain}
\DeviconCatalogEntry{qodana-plain-wordmark}
\DeviconCatalogEntry{qt-original}
\DeviconCatalogEntry{qt-plain}
\DeviconCatalogEntry{qtest-original}
\DeviconCatalogEntry{qtest-original-wordmark}
\DeviconCatalogEntry{qtest-plain}
\DeviconCatalogEntry{qtest-plain-wordmark}
\DeviconCatalogEntry{quarkus-plain}
\DeviconCatalogEntry{quarkus-plain-wordmark}
\DeviconCatalogEntry{quasar-plain}
\DeviconCatalogEntry{quasar-plain-wordmark}
\DeviconCatalogEntry{qwik-plain}
\DeviconCatalogEntry{qwik-plain-wordmark}
\DeviconCatalogEntry{r-line}
\DeviconCatalogEntry{r-plain}
\DeviconCatalogEntry{rabbitmq-original}
\DeviconCatalogEntry{rabbitmq-plain}
\DeviconCatalogEntry{rabbitmq-plain-wordmark}
\DeviconCatalogEntry{racket-line}
\DeviconCatalogEntry{racket-plain}
\DeviconCatalogEntry{radstudio-plain}
\DeviconCatalogEntry{rails-plain}
\DeviconCatalogEntry{rails-plain-wordmark}
\DeviconCatalogEntry{railway-line}
\DeviconCatalogEntry{railway-line-wordmark}
\DeviconCatalogEntry{railway-original}
\DeviconCatalogEntry{railway-original-wordmark}
\DeviconCatalogEntry{railway-plain}
\DeviconCatalogEntry{railway-plain-wordmark}
\DeviconCatalogEntry{rancher-line}
\DeviconCatalogEntry{rancher-line-wordmark}
\DeviconCatalogEntry{rancher-original}
\DeviconCatalogEntry{rancher-plain}
\DeviconCatalogEntry{rancher-plain-wordmark}
\DeviconCatalogEntry{raspberrypi-line}
\DeviconCatalogEntry{raspberrypi-line-wordmark}
\DeviconCatalogEntry{raspberrypi-plain}
\DeviconCatalogEntry{raspberrypi-plain-wordmark}
\DeviconCatalogEntry{ravendb-line}
\DeviconCatalogEntry{ravendb-line-wordmark}
\DeviconCatalogEntry{ravendb-original}
\DeviconCatalogEntry{ravendb-original-wordmark}
\DeviconCatalogEntry{ravendb-plain}
\DeviconCatalogEntry{ravendb-plain-wordmark}
\DeviconCatalogEntry{reach-plain}
\DeviconCatalogEntry{react-original}
\DeviconCatalogEntry{react-original-wordmark}
\DeviconCatalogEntry{react-plain}
\DeviconCatalogEntry{react-plain-wordmark}
\DeviconCatalogEntry{reactbootstrap-original}
\DeviconCatalogEntry{reactbootstrap-plain}
\DeviconCatalogEntry{reactnative-original}
\DeviconCatalogEntry{reactnative-original-wordmark}
\DeviconCatalogEntry{reactnavigation-original}
\DeviconCatalogEntry{reactnavigation-plain}
\DeviconCatalogEntry{reactrouter-plain}
\DeviconCatalogEntry{reactrouter-plain-wordmark}
\DeviconCatalogEntry{readthedocs-line}
\DeviconCatalogEntry{readthedocs-original}
\DeviconCatalogEntry{readthedocs-original-wordmark}
\DeviconCatalogEntry{readthedocs-plain}
\DeviconCatalogEntry{readthedocs-plain-wordmark}
\DeviconCatalogEntry{realm-plain}
\DeviconCatalogEntry{realm-plain-wordmark}
\DeviconCatalogEntry{rect-plain}
\DeviconCatalogEntry{redhat-plain}
\DeviconCatalogEntry{redhat-plain-wordmark}
\DeviconCatalogEntry{redis-plain}
\DeviconCatalogEntry{redis-plain-wordmark}
\DeviconCatalogEntry{redux-original}
\DeviconCatalogEntry{redux-plain}
\DeviconCatalogEntry{reflex-plain}
\DeviconCatalogEntry{reflex-plain-wordmark}
\DeviconCatalogEntry{remix-line}
\DeviconCatalogEntry{remix-line-wordmark}
\DeviconCatalogEntry{remix-original}
\DeviconCatalogEntry{remix-original-wordmark}
\DeviconCatalogEntry{remix-plain}
\DeviconCatalogEntry{remix-plain-wordmark}
\DeviconCatalogEntry{renpy-plain}
\DeviconCatalogEntry{replit-original}
\DeviconCatalogEntry{replit-plain}
\DeviconCatalogEntry{replit-plain-wordmark}
\DeviconCatalogEntry{rexx-plain}
\DeviconCatalogEntry{rexx-plain-wordmark}
\DeviconCatalogEntry{rider-plain}
\DeviconCatalogEntry{rider-plain-wordmark}
\DeviconCatalogEntry{rocksdb-line}
\DeviconCatalogEntry{rocksdb-plain}
\DeviconCatalogEntry{rockylinux-original}
\DeviconCatalogEntry{rockylinux-plain}
\DeviconCatalogEntry{rockylinux-plain-wordmark}
\DeviconCatalogEntry{rollup-line}
\DeviconCatalogEntry{rollup-line-wordmark}
\DeviconCatalogEntry{rollup-plain}
\DeviconCatalogEntry{rollup-plain-wordmark}
\DeviconCatalogEntry{ros-original}
\DeviconCatalogEntry{ros-original-wordmark}
\DeviconCatalogEntry{rspec-line}
\DeviconCatalogEntry{rspec-line-wordmark}
\DeviconCatalogEntry{rspec-plain}
\DeviconCatalogEntry{rspec-plain-wordmark}
\DeviconCatalogEntry{rstudio-plain}
\DeviconCatalogEntry{ruby-plain}
\DeviconCatalogEntry{ruby-plain-wordmark}
\DeviconCatalogEntry{rubymine-plain}
\DeviconCatalogEntry{rubymine-plain-wordmark}
\DeviconCatalogEntry{rust-line}
\DeviconCatalogEntry{rust-original}
\DeviconCatalogEntry{rust-plain}
\DeviconCatalogEntry{rxjs-plain}
\DeviconCatalogEntry{safari-line}
\DeviconCatalogEntry{safari-line-wordmark}
\DeviconCatalogEntry{safari-plain}
\DeviconCatalogEntry{safari-plain-wordmark}
\DeviconCatalogEntry{salesforce-plain}
\DeviconCatalogEntry{sanity-plain}
\DeviconCatalogEntry{sass-original}
\DeviconCatalogEntry{sass-plain}
\DeviconCatalogEntry{scala-plain}
\DeviconCatalogEntry{scala-plain-wordmark}
\DeviconCatalogEntry{scaleway-line}
\DeviconCatalogEntry{scaleway-line-wordmark}
\DeviconCatalogEntry{scaleway-original}
\DeviconCatalogEntry{scaleway-original-wordmark}
\DeviconCatalogEntry{scaleway-plain}
\DeviconCatalogEntry{scaleway-plain-wordmark}
\DeviconCatalogEntry{scalingo-line}
\DeviconCatalogEntry{scalingo-line-wordmark}
\DeviconCatalogEntry{scalingo-plain}
\DeviconCatalogEntry{scalingo-plain-wordmark}
\DeviconCatalogEntry{scikitlearn-line}
\DeviconCatalogEntry{scikitlearn-line-wordmark}
\DeviconCatalogEntry{scikitlearn-plain}
\DeviconCatalogEntry{scikitlearn-plain-wordmark}
\DeviconCatalogEntry{sdl-plain}
\DeviconCatalogEntry{seaborn-plain}
\DeviconCatalogEntry{seaborn-plain-wordmark}
\DeviconCatalogEntry{selenium-original}
\DeviconCatalogEntry{selenium-plain}
\DeviconCatalogEntry{sema-original}
\DeviconCatalogEntry{sema-original-wordmark}
\DeviconCatalogEntry{sema-plain}
\DeviconCatalogEntry{sema-plain-wordmark}
\DeviconCatalogEntry{sentry-original}
\DeviconCatalogEntry{sentry-original-wordmark}
\DeviconCatalogEntry{sentry-plain}
\DeviconCatalogEntry{sentry-plain-wordmark}
\DeviconCatalogEntry{sequelize-plain}
\DeviconCatalogEntry{sequelize-plain-wordmark}
\DeviconCatalogEntry{sfml-plain}
\DeviconCatalogEntry{shadcn-original}
\DeviconCatalogEntry{shopware-original}
\DeviconCatalogEntry{shopware-original-wordmark}
\DeviconCatalogEntry{shotgrid-original-wordmark}
\DeviconCatalogEntry{shotgrid-plain}
\DeviconCatalogEntry{shotgrid-plain-wordmark}
\DeviconCatalogEntry{sketch-line}
\DeviconCatalogEntry{sketch-line-wordmark}
\DeviconCatalogEntry{sketch-plain}
\DeviconCatalogEntry{sketch-plain-wordmark}
\DeviconCatalogEntry{slack-plain}
\DeviconCatalogEntry{slack-plain-wordmark}
\DeviconCatalogEntry{socketio-original}
\DeviconCatalogEntry{socketio-original-wordmark}
\DeviconCatalogEntry{socketio-plain}
\DeviconCatalogEntry{socketio-plain-wordmark}
\DeviconCatalogEntry{solidity-plain}
\DeviconCatalogEntry{solidjs-plain}
\DeviconCatalogEntry{solidjs-plain-wordmark}
\DeviconCatalogEntry{sonarqube-line}
\DeviconCatalogEntry{sonarqube-line-wordmark}
\DeviconCatalogEntry{sonarqube-original}
\DeviconCatalogEntry{sonarqube-plain}
\DeviconCatalogEntry{sonarqube-plain-wordmark}
\DeviconCatalogEntry{sourceengine-plain}
\DeviconCatalogEntry{sourceengine-plain-wordmark}
\DeviconCatalogEntry{sourcetree-original}
\DeviconCatalogEntry{sourcetree-original-wordmark}
\DeviconCatalogEntry{sourcetree-plain}
\DeviconCatalogEntry{sourcetree-plain-wordmark}
\DeviconCatalogEntry{spack-plain}
\DeviconCatalogEntry{spicedb-line}
\DeviconCatalogEntry{spicedb-plain}
\DeviconCatalogEntry{splunk-original-wordmark}
\DeviconCatalogEntry{splunk-plain-wordmark}
\DeviconCatalogEntry{spring-original}
\DeviconCatalogEntry{spring-original-wordmark}
\DeviconCatalogEntry{spring-plain}
\DeviconCatalogEntry{spring-plain-wordmark}
\DeviconCatalogEntry{spss-plain}
\DeviconCatalogEntry{spyder-plain}
\DeviconCatalogEntry{spyder-plain-wordmark}
\DeviconCatalogEntry{sqlalchemy-plain}
\DeviconCatalogEntry{sqlalchemy-plain-wordmark}
\DeviconCatalogEntry{sqldeveloper-plain}
\DeviconCatalogEntry{sqlite-plain}
\DeviconCatalogEntry{sqlite-plain-wordmark}
\DeviconCatalogEntry{ssh-original}
\DeviconCatalogEntry{ssh-original-wordmark}
\DeviconCatalogEntry{ssh-plain}
\DeviconCatalogEntry{ssh-plain-wordmark}
\DeviconCatalogEntry{stackblitz-line}
\DeviconCatalogEntry{stackblitz-line-wordmark}
\DeviconCatalogEntry{stackblitz-original}
\DeviconCatalogEntry{stackblitz-plain}
\DeviconCatalogEntry{stackblitz-plain-wordmark}
\DeviconCatalogEntry{stackoverflow-line}
\DeviconCatalogEntry{stackoverflow-line-wordmark}
\DeviconCatalogEntry{stackoverflow-plain}
\DeviconCatalogEntry{stackoverflow-plain-wordmark}
\DeviconCatalogEntry{starship-plain}
\DeviconCatalogEntry{starship-plain-wordmark}
\DeviconCatalogEntry{stata-original-wordmark}
\DeviconCatalogEntry{stata-plain-wordmark}
\DeviconCatalogEntry{stenciljs-plain}
\DeviconCatalogEntry{stenciljs-plain-wordmark}
\DeviconCatalogEntry{stimulus-original}
\DeviconCatalogEntry{stimulus-plain}
\DeviconCatalogEntry{stimulus-plain-wordmark}
\DeviconCatalogEntry{storybook-plain}
\DeviconCatalogEntry{storybook-plain-wordmark}
\DeviconCatalogEntry{streamlit-plain}
\DeviconCatalogEntry{streamlit-plain-wordmark}
\DeviconCatalogEntry{styledcomponents-plain}
\DeviconCatalogEntry{styledcomponents-plain-wordmark}
\DeviconCatalogEntry{stylus-original}
\DeviconCatalogEntry{subversion-original}
\DeviconCatalogEntry{subversion-plain}
\DeviconCatalogEntry{subversion-plain-wordmark}
\DeviconCatalogEntry{sulu-line}
\DeviconCatalogEntry{sulu-line-wordmark}
\DeviconCatalogEntry{sulu-original}
\DeviconCatalogEntry{sulu-original-wordmark}
\DeviconCatalogEntry{sulu-plain}
\DeviconCatalogEntry{sulu-plain-wordmark}
\DeviconCatalogEntry{supabase-plain}
\DeviconCatalogEntry{supabase-plain-wordmark}
\DeviconCatalogEntry{surrealdb-plain}
\DeviconCatalogEntry{surrealdb-plain-wordmark}
\DeviconCatalogEntry{svelte-plain}
\DeviconCatalogEntry{svelte-plain-wordmark}
\DeviconCatalogEntry{svgo-line}
\DeviconCatalogEntry{svgo-line-wordmark}
\DeviconCatalogEntry{svgo-plain}
\DeviconCatalogEntry{svgo-plain-wordmark}
\DeviconCatalogEntry{swagger-plain}
\DeviconCatalogEntry{swagger-plain-wordmark}
\DeviconCatalogEntry{swift-plain}
\DeviconCatalogEntry{swift-plain-wordmark}
\DeviconCatalogEntry{swiper-original}
\DeviconCatalogEntry{swiper-plain}
\DeviconCatalogEntry{symfony-original}
\DeviconCatalogEntry{symfony-original-wordmark}
\DeviconCatalogEntry{symfony-plain}
\DeviconCatalogEntry{symfony-plain-wordmark}
\DeviconCatalogEntry{systemverilog-original}
\DeviconCatalogEntry{systemverilog-plain}
\DeviconCatalogEntry{tailwindcss-original}
\DeviconCatalogEntry{tailwindcss-plain}
\DeviconCatalogEntry{tailwindcss-plain-wordmark}
\DeviconCatalogEntry{talos-plain}
\DeviconCatalogEntry{tauri-plain}
\DeviconCatalogEntry{tauri-plain-wordmark}
\DeviconCatalogEntry{teamcity-plain}
\DeviconCatalogEntry{teamcity-plain-wordmark}
\DeviconCatalogEntry{teleport-line}
\DeviconCatalogEntry{teleport-line-wordmark}
\DeviconCatalogEntry{teleport-original}
\DeviconCatalogEntry{teleport-original-wordmark}
\DeviconCatalogEntry{teleport-plain}
\DeviconCatalogEntry{teleport-plain-wordmark}
\DeviconCatalogEntry{tensorflow-line}
\DeviconCatalogEntry{tensorflow-line-wordmark}
\DeviconCatalogEntry{tensorflow-original}
\DeviconCatalogEntry{tensorflow-original-wordmark}
\DeviconCatalogEntry{tenzir-line}
\DeviconCatalogEntry{tenzir-line-wordmark}
\DeviconCatalogEntry{tenzir-original}
\DeviconCatalogEntry{tenzir-original-wordmark}
\DeviconCatalogEntry{tenzir-plain}
\DeviconCatalogEntry{tenzir-plain-wordmark}
\DeviconCatalogEntry{terraform-plain}
\DeviconCatalogEntry{terraform-plain-wordmark}
\DeviconCatalogEntry{terramate-original}
\DeviconCatalogEntry{terramate-original-wordmark}
\DeviconCatalogEntry{terramate-plain}
\DeviconCatalogEntry{terramate-plain-wordmark}
\DeviconCatalogEntry{testrail-original}
\DeviconCatalogEntry{testrail-plain}
\DeviconCatalogEntry{testrail-plain-wordmark}
\DeviconCatalogEntry{tex-original}
\DeviconCatalogEntry{tex-original-wordmark}
\DeviconCatalogEntry{tex-plain}
\DeviconCatalogEntry{tex-plain-wordmark}
\DeviconCatalogEntry{thealgorithms-plain}
\DeviconCatalogEntry{thealgorithms-plain-wordmark}
\DeviconCatalogEntry{threedsmax-plain}
\DeviconCatalogEntry{threedsmax-plain-wordmark}
\DeviconCatalogEntry{threejs-original}
\DeviconCatalogEntry{threejs-original-wordmark}
\DeviconCatalogEntry{thymeleaf-plain}
\DeviconCatalogEntry{thymeleaf-plain-wordmark}
\DeviconCatalogEntry{tigerbeetle-line}
\DeviconCatalogEntry{tigerbeetle-original}
\DeviconCatalogEntry{tigerbeetle-original-wordmark}
\DeviconCatalogEntry{tigerbeetle-plain}
\DeviconCatalogEntry{tigerbeetle-plain-wordmark}
\DeviconCatalogEntry{tinygo-line}
\DeviconCatalogEntry{tinygo-plain}
\DeviconCatalogEntry{titaniumsdk-original}
\DeviconCatalogEntry{titaniumsdk-plain}
\DeviconCatalogEntry{tmux-plain}
\DeviconCatalogEntry{tmux-plain-wordmark}
\DeviconCatalogEntry{tomcat-line}
\DeviconCatalogEntry{tomcat-line-wordmark}
\DeviconCatalogEntry{toml-plain}
\DeviconCatalogEntry{tortoisegit-line}
\DeviconCatalogEntry{tortoisegit-plain}
\DeviconCatalogEntry{towergit-plain}
\DeviconCatalogEntry{towergit-plain-wordmark}
\DeviconCatalogEntry{traefikmesh-line}
\DeviconCatalogEntry{traefikmesh-line-wordmark}
\DeviconCatalogEntry{traefikmesh-original}
\DeviconCatalogEntry{traefikmesh-plain}
\DeviconCatalogEntry{traefikmesh-plain-wordmark}
\DeviconCatalogEntry{traefikproxy-line}
\DeviconCatalogEntry{traefikproxy-line-wordmark}
\DeviconCatalogEntry{traefikproxy-original}
\DeviconCatalogEntry{traefikproxy-plain}
\DeviconCatalogEntry{traefikproxy-plain-wordmark}
\DeviconCatalogEntry{travis-line}
\DeviconCatalogEntry{travis-line-wordmark}
\DeviconCatalogEntry{travis-plain}
\DeviconCatalogEntry{travis-plain-wordmark}
\DeviconCatalogEntry{trello-line}
\DeviconCatalogEntry{trello-line-wordmark}
\DeviconCatalogEntry{trello-plain}
\DeviconCatalogEntry{trello-plain-wordmark}
\DeviconCatalogEntry{trpc-plain}
\DeviconCatalogEntry{trpc-plain-wordmark}
\DeviconCatalogEntry{turbo-original}
\DeviconCatalogEntry{turbo-plain}
\DeviconCatalogEntry{turbo-plain-wordmark}
\DeviconCatalogEntry{twilio-original}
\DeviconCatalogEntry{twilio-original-wordmark}
\DeviconCatalogEntry{twilio-plain}
\DeviconCatalogEntry{twilio-plain-wordmark}
\DeviconCatalogEntry{twitter-original}
\DeviconCatalogEntry{twitter-plain}
\DeviconCatalogEntry{typeorm-line}
\DeviconCatalogEntry{typeorm-line-wordmark}
\DeviconCatalogEntry{typeorm-plain}
\DeviconCatalogEntry{typeorm-plain-wordmark}
\DeviconCatalogEntry{typer-line}
\DeviconCatalogEntry{typer-line-wordmark}
\DeviconCatalogEntry{typer-plain}
\DeviconCatalogEntry{typer-plain-wordmark}
\DeviconCatalogEntry{typescript-original}
\DeviconCatalogEntry{typescript-plain}
\DeviconCatalogEntry{typo3-line}
\DeviconCatalogEntry{typo3-line-wordmark}
\DeviconCatalogEntry{typo3-original}
\DeviconCatalogEntry{typo3-plain}
\DeviconCatalogEntry{typo3-plain-wordmark}
\DeviconCatalogEntry{ubuntu-plain}
\DeviconCatalogEntry{ubuntu-plain-wordmark}
\DeviconCatalogEntry{unifiedmodelinglanguage-plain}
\DeviconCatalogEntry{unifiedmodelinglanguage-plain-wordmark}
\DeviconCatalogEntry{unity-line}
\DeviconCatalogEntry{unity-line-wordmark}
\DeviconCatalogEntry{unity-plain}
\DeviconCatalogEntry{unity-plain-wordmark}
\DeviconCatalogEntry{unix-original}
\DeviconCatalogEntry{unrealengine-original}
\DeviconCatalogEntry{unrealengine-original-wordmark}
\DeviconCatalogEntry{unrealengine-plain}
\DeviconCatalogEntry{unrealengine-plain-wordmark}
\DeviconCatalogEntry{uv-line}
\DeviconCatalogEntry{uv-original}
\DeviconCatalogEntry{uv-plain}
\DeviconCatalogEntry{uwsgi-plain}
\DeviconCatalogEntry{v8-plain}
\DeviconCatalogEntry{v8-plain-wordmark}
\DeviconCatalogEntry{vaadin-original}
\DeviconCatalogEntry{vaadin-plain}
\DeviconCatalogEntry{vagrant-plain}
\DeviconCatalogEntry{vagrant-plain-wordmark}
\DeviconCatalogEntry{vala-plain}
\DeviconCatalogEntry{vala-plain-wordmark}
\DeviconCatalogEntry{vapor-plain}
\DeviconCatalogEntry{vapor-plain-wordmark}
\DeviconCatalogEntry{vault-original}
\DeviconCatalogEntry{vault-plain}
\DeviconCatalogEntry{vault-plain-wordmark}
\DeviconCatalogEntry{veevalidate-line}
\DeviconCatalogEntry{veevalidate-original}
\DeviconCatalogEntry{veevalidate-plain}
\DeviconCatalogEntry{vercel-line}
\DeviconCatalogEntry{vercel-line-wordmark}
\DeviconCatalogEntry{vercel-original}
\DeviconCatalogEntry{vercel-original-wordmark}
\DeviconCatalogEntry{vercel-plain}
\DeviconCatalogEntry{vercel-plain-wordmark}
\DeviconCatalogEntry{verilog-original}
\DeviconCatalogEntry{verilog-plain}
\DeviconCatalogEntry{vertx-line}
\DeviconCatalogEntry{vertx-line-wordmark}
\DeviconCatalogEntry{vertx-plain}
\DeviconCatalogEntry{vertx-plain-wordmark}
\DeviconCatalogEntry{vim-plain}
\DeviconCatalogEntry{visualbasic-line}
\DeviconCatalogEntry{visualbasic-plain}
\DeviconCatalogEntry{visualstudio-line}
\DeviconCatalogEntry{visualstudio-plain}
\DeviconCatalogEntry{vite-original}
\DeviconCatalogEntry{vite-original-wordmark}
\DeviconCatalogEntry{vite-plain}
\DeviconCatalogEntry{vite-plain-wordmark}
\DeviconCatalogEntry{vitejs-plain}
\DeviconCatalogEntry{vitess-plain}
\DeviconCatalogEntry{vitess-plain-wordmark}
\DeviconCatalogEntry{vitest-plain}
\DeviconCatalogEntry{vlang-original}
\DeviconCatalogEntry{vlang-plain}
\DeviconCatalogEntry{vscode-plain}
\DeviconCatalogEntry{vscode-plain-wordmark}
\DeviconCatalogEntry{vscodium-plain}
\DeviconCatalogEntry{vsphere-line}
\DeviconCatalogEntry{vsphere-line-wordmark}
\DeviconCatalogEntry{vsphere-plain}
\DeviconCatalogEntry{vsphere-plain-wordmark}
\DeviconCatalogEntry{vuejs-line}
\DeviconCatalogEntry{vuejs-line-wordmark}
\DeviconCatalogEntry{vuejs-plain}
\DeviconCatalogEntry{vuejs-plain-wordmark}
\DeviconCatalogEntry{vuestorefront-plain}
\DeviconCatalogEntry{vuetify-line}
\DeviconCatalogEntry{vuetify-plain}
\DeviconCatalogEntry{vulkan-line}
\DeviconCatalogEntry{vulkan-line-wordmark}
\DeviconCatalogEntry{vulkan-original}
\DeviconCatalogEntry{vulkan-original-wordmark}
\DeviconCatalogEntry{vulkan-plain}
\DeviconCatalogEntry{vulkan-plain-wordmark}
\DeviconCatalogEntry{vyper-original}
\DeviconCatalogEntry{vyper-original-wordmark}
\DeviconCatalogEntry{vyper-plain}
\DeviconCatalogEntry{vyper-plain-wordmark}
\DeviconCatalogEntry{waku-line}
\DeviconCatalogEntry{waku-plain}
\DeviconCatalogEntry{wasm-original}
\DeviconCatalogEntry{wasm-plain}
\DeviconCatalogEntry{wasm-plain-wordmark}
\DeviconCatalogEntry{wasp-plain}
\DeviconCatalogEntry{web2py-plain}
\DeviconCatalogEntry{web3js-plain}
\DeviconCatalogEntry{webflow-original}
\DeviconCatalogEntry{webgl-original}
\DeviconCatalogEntry{webgl-plain}
\DeviconCatalogEntry{webgpu-line}
\DeviconCatalogEntry{webgpu-line-wordmark}
\DeviconCatalogEntry{webgpu-plain}
\DeviconCatalogEntry{webgpu-plain-wordmark}
\DeviconCatalogEntry{weblate-plain}
\DeviconCatalogEntry{weblate-plain-wordmark}
\DeviconCatalogEntry{webpack-plain}
\DeviconCatalogEntry{webpack-plain-wordmark}
\DeviconCatalogEntry{webstorm-plain}
\DeviconCatalogEntry{webstorm-plain-wordmark}
\DeviconCatalogEntry{windows11-original}
\DeviconCatalogEntry{windows11-original-wordmark}
\DeviconCatalogEntry{windows11-plain}
\DeviconCatalogEntry{windows11-plain-wordmark}
\DeviconCatalogEntry{windows8-original}
\DeviconCatalogEntry{windows8-original-wordmark}
\DeviconCatalogEntry{windows8-plain}
\DeviconCatalogEntry{windows8-plain-wordmark}
\DeviconCatalogEntry{wireguard-line}
\DeviconCatalogEntry{wireguard-original}
\DeviconCatalogEntry{wireguard-plain}
\DeviconCatalogEntry{wireguard-plain-wordmark}
\DeviconCatalogEntry{wolfram-plain}
\DeviconCatalogEntry{wolfram-plain-wordmark}
\DeviconCatalogEntry{woocommerce-plain}
\DeviconCatalogEntry{woocommerce-plain-wordmark}
\DeviconCatalogEntry{wordpress-plain}
\DeviconCatalogEntry{wordpress-plain-wordmark}
\DeviconCatalogEntry{xamarin-original}
\DeviconCatalogEntry{xamarin-original-wordmark}
\DeviconCatalogEntry{xamarin-plain}
\DeviconCatalogEntry{xamarin-plain-wordmark}
\DeviconCatalogEntry{xcode-plain}
\DeviconCatalogEntry{xd-line}
\DeviconCatalogEntry{xd-plain}
\DeviconCatalogEntry{xml-line}
\DeviconCatalogEntry{xml-line-wordmark}
\DeviconCatalogEntry{xml-plain}
\DeviconCatalogEntry{xml-plain-wordmark}
\DeviconCatalogEntry{yaml-plain}
\DeviconCatalogEntry{yaml-plain-wordmark}
\DeviconCatalogEntry{yarn-line}
\DeviconCatalogEntry{yarn-line-wordmark}
\DeviconCatalogEntry{yarn-original}
\DeviconCatalogEntry{yarn-original-wordmark}
\DeviconCatalogEntry{yarn-plain}
\DeviconCatalogEntry{yarn-plain-wordmark}
\DeviconCatalogEntry{yii-plain}
\DeviconCatalogEntry{yii-plain-wordmark}
\DeviconCatalogEntry{youtrack-plain}
\DeviconCatalogEntry{youtrack-plain-wordmark}
\DeviconCatalogEntry{yugabytedb-plain}
\DeviconCatalogEntry{yugabytedb-plain-wordmark}
\DeviconCatalogEntry{yunohost-plain}
\DeviconCatalogEntry{zed-original}
\DeviconCatalogEntry{zed-plain}
\DeviconCatalogEntry{zend-line}
\DeviconCatalogEntry{zend-line-wordmark}
\DeviconCatalogEntry{zend-original}
\DeviconCatalogEntry{zend-original-wordmark}
\DeviconCatalogEntry{zend-plain}
\DeviconCatalogEntry{zend-plain-wordmark}
\DeviconCatalogEntry{zig-original}
\DeviconCatalogEntry{zig-plain}
\DeviconCatalogEntry{zig-plain-wordmark}
\DeviconCatalogEntry{zsh-line}
\DeviconCatalogEntry{zsh-line-wordmark}
\DeviconCatalogEntry{zsh-plain}
\DeviconCatalogEntry{zsh-plain-wordmark}
\DeviconCatalogEntry{zustand-plain}

\end{multicols}
\end{document}
